\chapter{Introdução}

Os filtros passivos, como as topologias RC, são blocos fundamentais em sistemas eletrônicos, sendo amplamente empregados no condicionamento de sinais e na atenuação de componentes de alta frequência indesejadas. A análise destes circuitos no domínio da frequência, frequentemente facilitada pelo uso da Transformada de Laplace, permite prever o seu comportamento dinâmico e determinar parâmetros cruciais, como a frequência de corte e o ganho estático.

Contudo, a aplicação prática de filtros passivos apresenta desafios significativos quando estes são acoplados a outros estágios de um circuito. A conexão de uma carga resistiva direta altera a impedância equivalente do sistema, provocando o desvio da frequência de corte original e a atenuação do ganho, um fenômeno conhecido como efeito de carga. Para mitigar este problema, recorre-se frequentemente ao uso de amplificadores operacionais configurados como seguidores de tensão (\textit{buffers}). Devido à sua elevada impedância de entrada e baixa impedância de saída, o \textit{buffer} atua como um isolador ideal, garantindo que o filtro opere de forma independente da carga conectada.

Nesta prática laboratorial, realizou-se o projeto, a análise matemática teórica e a validação experimental de um filtro passa-baixa RC nas suas configurações sem carga, com carga resistiva direta e com isolamento ativo via \textit{buffer}. O estudo teórico foi suportado pela ferramenta de computação simbólica Maxima \cite{maxima} para a resolução algébrica e determinação das funções de transferência. Posteriormente, os ensaios de bancada permitiram extrair o diagrama de Bode de magnitude real do circuito e confrontá-lo com o modelo analítico.

\section{Objetivos}

A presente prática teve como objetivo central explorar as características operacionais e os limites de aplicação de filtros passivos de primeira ordem. Os objetivos específicos definidos para o trabalho incluem:

\begin{itemize}
	\item Compreender e caracterizar o comportamento dinâmico de um circuito RC operando como um filtro passa-baixa;
	\item Aplicar a Transformada de Laplace como ferramenta principal na análise de circuitos e na obtenção das funções de transferência;
	\item Utilizar o software de computação simbólica Maxima \cite{maxima} para resolver as equações algébricas resultantes da modelagem linear do circuito;
	\item Treinar a análise de circuitos elétricos submetidos a condições de carga e alimentados por fontes não ideais;
	\item Empregar um amplificador operacional na configuração de \textit{buffer} para minimizar o impacto de uma carga no comportamento do filtro;
	\item Realizar medições experimentais das grandezas elétricas utilizando um osciloscópio e um gerador de sinais;
	\item Levantar o comportamento na frequência dos circuitos de filtragem em bancada para comparar e validar a magnitude de resposta em frequência esperada de forma teórica.
\end{itemize}